\chapter{Bases, dimensión y coordenadas}
\label{cap:bases-dimension}

Las bases proporcionan descripciones eficientes de los vectores y permiten asignar una dimensión finita al espacio. Se formalizan la extensión de conjuntos independientes, el intercambio y el cambio de coordenadas \cite{lang1986,lang1987,axler2024}.

\section{Bases de un espacio vectorial}
Conjuntos que generan el espacio y son linealmente independientes.

\section{Existencia y extensión de bases}
Construcción de bases en dimensión finita y extensión de conjuntos independientes.

\section{Lema de intercambio}
Intercambio de generadores por vectores independientes y sus consecuencias.

\section{Dimensión}
Invariancia del número de elementos de una base y dimensión finita.

\section{Dimensión de subespacios}
Comparación de dimensiones, inclusión y fórmulas para sumas e intersecciones.

\section{Coordenadas respecto de una base}
Vector de coordenadas, unicidad y aplicaciones a cálculos.

\section{Matrices de cambio de coordenadas}
Conversión entre sistemas de coordenadas y composición de cambios.
